\chapter{Threat Model and Problem Statement}\label{ch:threat}

This chapter defines the actors, adversaries, assumptions, security goals, and scope that govern the
guarantees evaluated in later chapters. It combines the malicious-directory model of key transparency
with ACKA's active-network model so that the design can be evaluated against both server-side and
in-transit key substitution.

% ============================================================================
\section{System Model and Actors}\label{sec:threat-system}

DUET considers asynchronous communication between Alice and Bob over
infrastructure operated by a service provider (\cref{sec:bg-signal}). Either party
may be offline while the other sends a message or updates key material. The system
model comprises four actors: the
communicating clients, the messaging infrastructure, the key-transparency
service, and independent auditors. Their respective responsibilities and trust
assumptions define the attack surface considered throughout the remainder of this
chapter.

The \textbf{clients}, Alice ($A$) and Bob ($B$), run the messaging application on
their own devices. They generate and hold their own secret keys, publish the
corresponding public key material, and perform all verification locally. They are
assumed honest: their devices are not under adversarial control (the boundary
where this fails is scenario X1 in \cref{sec:threat-adversary}). Clients may be
resource-constrained and intermittently connected, which is why the cost of
catching up after an offline period is a primary concern.

The \textbf{message-delivery and key-distribution server} is the provider's
infrastructure. It stores and serves the X3DH pre-key bundles that Alice needs to
start a session with Bob (\cref{ssec:bg-x3dh}) and relays ciphertext between the
parties. It is the single service-side point through which public keys reach their consumers,
and therefore the single service-side point at which a key can be substituted. The
client--service channel is the other substitution surface, modelled separately as
the network adversary (A3, \cref{sec:threat-adversary}).

The \textbf{key-transparency log server} maintains the verifiable directory and
its append-only log, answers lookup and monitoring queries with cryptographic
proofs, and publishes a signed tree head each epoch. It may be operated by the
same party as the message-delivery server. The analysis assumes the worst case
throughout: that the two are a single, possibly malicious, operator, which avoids
relying on any assumption that they do not collude.

The \textbf{auditors} are independent parties that observe the log from a second
vantage point, cross-check the tree heads the server publishes, and replay each
epoch's declared update set against the committed directory root (the per-epoch
transition audit of \cref{ch:design}). Their role in this model is fork detection
(scenario A4) and binding a malicious operator to declare every cross-epoch
mutation. The dual-index warning itself depends on neither, a separation made
precise in \cref{sec:threat-assumptions}.

The interactions between these actors define three trust boundaries
(\cref{fig:threat-arch}). The
\emph{client--server} boundary is semi-trusted: the standard key-transparency
trust configuration inherited from CONIKS and the IETF architecture (\cref{sec:bg-logs}), in
which the server is relied upon for availability but not for honesty, and not the
abstract, security-defined log that ACKA leaves to be implemented. The \emph{client--network} boundary is
untrusted: the full active network adversary described below is adopted, the same
adversary ACKA defends against in its continuous mode (\cref{sec:bg-acka}). The
\emph{server--auditor} boundary is independent: the auditor's value is that its
view of the log is obtained separately from the client's.

The system operates in one of two authentication modes, both developed in
\cref{ch:design} and referred to in this chapter only by name: a \emph{session
mode}, in which a per-conversation dual-index entry authenticates the keys of a
given session, and a \emph{continuous mode}, in which a per-ratchet-epoch
fingerprint authenticates the evolving Double Ratchet state
(\cref{ssec:bg-ratchet}). The threat model applies to both. Where a guarantee is
specific to one mode, this is indicated.

\begin{figure}[t]
\centering
% fig-threat-arch.tex — trust boundaries.
% The untrusted network is a band BETWEEN the honest clients and the operator;
% the operator sits wholly outside it, and the auditor is beyond the operator.
% Label: \label{fig:threat-arch}
\begin{tikzpicture}[>=Stealth,font=\small,
   actor/.style={draw,rounded corners,minimum width=2.3cm,minimum height=10mm,align=center}]
  \node[actor] (A) at (0,0.95)  {Alice\\\footnotesize(client)};
  \node[actor] (B) at (0,-0.95) {Bob\\\footnotesize(client)};
  \node[actor,fill=black!8] (S) at (6.0,0) {Service operator\\\footnotesize message delivery\\\footnotesize {}+ KT log server};
  \node[actor] (Au) at (10.6,0) {Auditor};
  % untrusted-network band: the medium between clients and operator only, not the
  % honest endpoints; operator wholly to its right.
  \fill[red!4,rounded corners] (1.55,-1.7) rectangle (3.75,1.7);
  \draw[dashed,red!70!black,rounded corners] (1.55,-1.7) rectangle (3.75,1.7);
  \node[red!70!black,font=\footnotesize,align=center] at (2.65,2.05) {untrusted network\\(Dolev--Yao)};
  % client--server links pass through the band
  \draw[->] (A.east) -- (S.west);
  \draw[->] (B.east) -- (S.west);
  \node[font=\footnotesize,align=center] at (6.0,-1.9) {service operator:\\semi-trusted};
  % independent auditor beyond the operator boundary
  \draw[<->,dashed,green!45!black] (S.east) -- node[above,font=\footnotesize]{independent} (Au.west);
\end{tikzpicture}

\caption[System model and trust boundaries.]{System model and trust boundaries. Clients reach the service operator across a
network controlled by a Dolev--Yao adversary. The operator provides message delivery
and the key-transparency log and is treated as one semi-trusted party. An independent
auditor observes the log from a separate vantage point.}
\label{fig:threat-arch}
\end{figure}

% ============================================================================
\FloatBarrier
\section{Adversary Model}\label{sec:threat-adversary}

The threat model combines two well-established adversary models, because the
system must reason about both network attacks and malicious key-directory
behaviour. For the network, the analysis adopts the Dolev--Yao
model~\cite{dolevyao1983}. For the infrastructure, it adopts the
semi-trusted-server model of the key-transparency literature. These adversaries
are modelled separately, and deployments may face either or both simultaneously.

\paragraph{The Network Attacker (Dolev--Yao).} Following Dolev and
Yao~\cite{dolevyao1983}, the network adversary has maximal control of the wire:
it can intercept, modify, inject, delete, and replay any message, and it can
impersonate any party in transit. It does not, on its own, control the server's
directory. This models the active man-in-the-middle that the Double Ratchet's
post-compromise security does not defend against once the adversary remains active
(\cref{ssec:bg-ratchet}), and that ACKA's continuous mode is designed to catch
(\cref{sec:bg-acka}).

\paragraph{The Infrastructure Attacker (Malicious Server).} The server is
semi-trusted: relied upon for availability but assumed to act maliciously when it
can do so without being caught. It may substitute the key it serves for a user,
censor writes a client attempts to make to the log, and equivocate by presenting
two internally consistent but divergent views of the directory to different
clients. This is the adversary key transparency was built to expose
(\cref{sec:bg-vkd}), and it lies outside ACKA's detection theorem, which assumes a $\Logsec$-secure log rather than constructing one against such a server.

\paragraph{What the Adversary Cannot Do.} The model grants the adversary the
capabilities above and no more. It cannot break the underlying primitives: it
cannot find collisions in SHA-256, solve the discrete-logarithm problem
underlying X25519, or forge an Ed25519 signature. It cannot forge a
verifiable-random-function proof without the server's VRF secret key, nor produce
a Merkle inclusion proof for an entry that was never committed to the log. It
cannot sustain equivocation without detection by an honest auditor holding a valid,
independent second-vantage head (\cref{sec:threat-goals}). And it cannot
compromise an honest client's device: endpoint security is out of scope, and the
boundary case where this assumption is lifted is scenario X1.

The implications of these adversary capabilities for the design are developed in
\cref{ch:design}.

\Cref{tab:threat-attacks} summarizes the attack scenarios this adversary model
admits, ranging from the honest baseline (A0) through the
attacks the system is designed to detect (A1--A5) to the boundary case it does not
address (X1). Each scenario is evaluated experimentally in \cref{ch:eval}.

\begin{table}[H]
\centering
\caption{Attack scenarios, adversaries, and detection mechanisms.}
\label{tab:threat-attacks}
\setlength{\extrarowheight}{2pt}
\small
\begin{tabularx}{\textwidth}{@{}ll>{\raggedright\arraybackslash}X>{\raggedright\arraybackslash}X@{}}
\toprule
Scenario & Adversary & Objective & Detection mechanism \\
\midrule
A0 & none & honest baseline & all checks pass \\
A1 & malicious server & substitute a user's key & self-check + peer warning \\
A2 & malicious server & substitute a key and censor the warning & substitution + missing warning \\
A3 & network (Dolev--Yao) & substitute a key in transit & divergent labels + global identity binding \\
A4 & malicious server & equivocate (split-view) & cross-vantage auditor \\
A5 & malicious server & substitute during a key rotation & identity re-fetch + predecessor cross-signature + warning \\
X1 & endpoint holder & impersonate using the victim's key & out of scope (boundary) \\
\bottomrule
\end{tabularx}
\end{table}

\FloatBarrier
\subsection{Attack Scenarios}\label{ssec:threat-scenarios}

The scenarios below turn the adversary model of \cref{tab:threat-attacks} into concrete attacks,
ordered from the primary key-substitution threat through its network and equivocation variants to the
endpoint-compromise boundary.

\paragraph{A1: Key Substitution.} A1 is the primary attack in the model. A
compromised or coerced server writes its own key in place of the victim's. The
substituted key is internally consistent: the malicious server has
legitimately appended the forged key to its own Merkle log and can produce valid
inclusion proofs for it, so the Merkle proofs alone do not detect A1. Detection instead comes from
the victim observing the wrong key under its own name and from the peer learning
of the substitution through the in-band warning channel. The evaluation
(\cref{ch:eval}) instantiates A1 in two forms: A1a, substitution of the global
identity slot, and A1b, substitution of an immutable dual-index session slot.

\paragraph{A2: Warning Censorship.} A2 concerns censorship before the warning commits. After
detecting a key mismatch, the client verifies the warning slot. A proof-verified empty slot raises the
compound A2 alarm. If a committed warning later disappears, the transition audit detects the invalid
directory transition. For a client using windowed catch-up, an in-window substitution is detected only
if the warning commits and remains present at the catch-up read (\cref{ssec:design-catchup}).

\paragraph{A3: Network Man-in-the-Middle.} With no server collusion, a
Dolev--Yao adversary substitutes a key in the bundle in transit. The two parties
then derive different shared secrets, and because the dual-index labels are
derived from that secret, they address disjoint locations in the directory: the
victim queries a label at which the impersonator never registered. In addition, the substituted
bundle diverges from the identity entry at the victim's global VRF-derived label, so the client's
comparison against that VRF-bound entry exposes the substitution independently of the missed
dual-index slot.

A3 also has a post-compromise variant that motivates the continuous mode: an
adversary that has already compromised a session and then injects a forged ratchet
key mid-conversation is invisible to a server-side audit, because the server and
log remain honest. Continuous mode detects it because each epoch logs a
fingerprint of the ratchet key, so an injected key produces a conflicting
fingerprint at the agreed label and the same warning is raised.

\paragraph{A4: Equivocation (Split-View).} Inherited from Certificate
Transparency (\cref{ssec:bg-ct2kt}), the server signs two heads at one epoch, one
for the victim and one for the auditor, and no single client can distinguish the
two views. The cross-vantage comparison of the two heads exposes the fork.

\paragraph{A5: Malicious Key Rotation.} During a legitimate key rotation,
changes to the directory are expected. A malicious server exploits this by
substituting its own key while presenting the rotation as genuine, so the peer
derives the wrong secret and lands at a different directory index. The defence
anchors a rotation to identity rather than to the session: on a rotation the peer
re-fetches the current key from the global VRF directory by user identity and compares
it against the key it already trusts, so a substituted rotation is detected as a changed
identity key rather than being accepted silently, and the dual-index warning is
the in-band backstop. The prototype's client relies on the identity re-fetch and the warning, and a
single-link cross-signature check, the new key signed by its trusted predecessor, is demonstrated in
the adversarial suite.

\paragraph{X1: Endpoint Compromise (Boundary).} Once the adversary holds the
victim's own private key, it derives the same shared secret, the same labels, and
the same proofs, and is indistinguishable from the victim. X1 marks the boundary of
key-transparency protection. Key transparency protects communication between honest
endpoints but cannot protect a compromised endpoint. The attack--defence structure of these
scenarios, separating the interposition routes that place an attacker key (A1, A3,
A5) from the concealment attempts that try to hide it (A2, A4) and from the
out-of-scope endpoint boundary (X1), is presented with the evaluation in
\cref{fig:eval-attacktree}.

% ============================================================================
\section{Justification of the Threat Model}\label{sec:threat-fair}

\paragraph{The Malicious-Server Assumption in Practice.} This is the standard
assumption of the key-transparency field. The same semi-trusted, potentially
equivocating server appears in CONIKS~\cite{melara2015coniks},
SEEMless~\cite{chase2019seemless}, Parakeet~\cite{malvai2023parakeet}, and the
IETF Key Transparency architecture, and the active network adversary is the one
ACKA~\cite{dowling2023acka} is designed to defend against. The assumption also
reflects deployment practice. Deployed end-to-end encrypted messengers commonly
rely on a key-distribution service to return users' public keys, and compromise or
coercion of that service can enable key substitution. Documented fake-key attacks
show the threat is practical rather than
hypothetical~\cite{yadav2022fakekey}. Assuming the server may misbehave is therefore a
conservative and well-precedented modelling choice.

\paragraph{Rejection of the Honest-Routing Assumption.} The model allows a network adversary to drop
messages, since no cryptographic directory can force their delivery. DUET makes such interference
observable. Censored writes, missing epochs, and invalid proofs are reported as failed checks rather
than successful verification.

One limitation remains, inherited from the trust-on-first-use assumption that also
underlies ACKA. Detection establishes that the key seen now differs from a key
trusted earlier. It cannot vouch for the first key, before any trusted baseline
exists. A substitution that predates first contact, poisoning the initial pin, is
not caught, because there is nothing to compare against. This limitation is
carried through the analysis and revisited in the discussion of limitations
(\cref{ch:discussion}).

% ============================================================================
\section{Assumptions}\label{sec:threat-assumptions}

The guarantees of the later chapters rest on the following assumptions.

The \textbf{cryptographic assumptions} are conventional: SHA-256 is
collision-resistant, so that a published commitment binds its contents
(\cref{ssec:bg-merkle}). Ed25519 is existentially unforgeable under chosen-message
attack, so that a signed tree head cannot be forged~\cite{rfc8032}, and the
verifiable random function provides both uniqueness and pseudorandomness as
specified in RFC~9381~\cite{rfc9381}, so that directory labels are deterministic
for the operator yet unpredictable to outsiders. The append-only consistency
property guaranteed by RFC~9162 consistency proofs~\cite{rfc9162} ensures that a
verified log cannot be silently rewritten between visits. These assumptions
underlie the adversary's capability limits in \cref{sec:threat-adversary}.

The \textbf{endpoint assumption} is that honest clients' devices are not
compromised. This is the X1 boundary, and lifting it places the situation outside
the scope of any key-transparency guarantee.

The \textbf{auditor assumption} is that at least one honest auditor observes the log. It is required
only for fork detection (A4) and the per-epoch transition audit (\cref{ch:design}), not for the
dual-index warning mechanism. Detection of server-side substitution (A1) and network
man-in-the-middle (A3) is performed by the communicating clients and depends on no auditor.

Finally, the model follows a trust-on-first-use (TOFU) approach. Peer-side detection acts only after a
key has been pinned. The resulting first-key limitation is justified in \cref{sec:threat-fair} and
revisited in \cref{ch:discussion}.

% ============================================================================
\FloatBarrier
\section{Security Goals}\label{sec:threat-goals}

The model defines four security goals, formalized in \cref{app:definitions} and specified in
\cref{ch:design}. \Cref{ch:eval} evaluates G1 and G2 experimentally. G3 rests on VRF pseudorandomness
with RFC~9381 conformance evidence, while G4 is evaluated under the second-vantage condition.

\paragraph{G1: Key-Substitution Detection.} If a server (A1, A5) or a network
adversary (A3) substitutes a key, an honest party detects it once an honest monitor completes the
required verified lookups and, where the peer's warning is needed, that warning write commits. For the established-pair attacks A1 and A5, G1 provides bidirectional detection through the in-band
warning. The victim detects the wrong key under its own name, and a peer that has pinned the
victim's key detects it on its next key-fetch. Either can raise the shared warning, so an offline
victim remains covered. First-contact substitution (A3) is instead caught by the divergence of the
secret-derived lookup label together with comparison against the global VRF-bound identity entry, not
by the warning channel. The scripted demonstration flags the change within one or two epochs
of the next check, which is a demonstration schedule rather than a protocol-wide latency bound.

\paragraph{G2: Consistency Across Offline Windows.} A client returning after an
offline period requests one distant-head consistency proof rather than one per missed epoch, and
confirms that the log it now sees is an append-only extension of the one it last verified. The
proof's size depends on the endpoint log sizes rather than on a proof count fixed directly to the
window $W$. The endpoints themselves grow as the log grows. This property ensures that intermittent connectivity does not impose a
verification burden that grows with the number of missed epochs. The residual growth is the
logarithmic dependence on the endpoint log sizes.

\paragraph{G3: Directory Privacy.} The mapping between a user's identity and its
directory slot is hidden from outsiders, via VRF-derived labels
(\cref{ssec:bg-ct2kt}). The limit of this goal is that it provides
directory privacy against third parties, not query privacy against the operator,
which necessarily learns which identity it is asked to resolve. Query privacy from
the operator is out of scope (\cref{sec:threat-scope}).

\paragraph{G4: Fork Detection.} An equivocating server (A4) that signs two
divergent heads at one epoch is detected by cross-vantage auditing when a valid, independent second-vantage head is available.
Absent one, the auditor reports \textsc{ConsistentNoVantage} rather than a fork verdict
(\cref{ch:impl}). G4 is the only
goal that depends on the auditor-honesty assumption (\cref{sec:threat-assumptions}).

\begin{table}[H]
\centering
\caption{Security goals, the attacks they address, and the mechanisms that deliver
them.}
\label{tab:threat-goals}
\setlength{\extrarowheight}{2pt}
\footnotesize
\begin{tabularx}{\textwidth}{@{}ll>{\raggedright\arraybackslash}Xl@{}}
\toprule
Goal & Attack(s) & Security mechanism & Auditor? \\
\midrule
G1 key-substitution detection & A1, A3, A5 & A1 and A5 use self-checks and the dual-index warning.\newline A3 uses divergent labels and the global VRF-bound identity entry.\newline A5 also uses identity re-fetch and predecessor cross-signature. & no \\
G2 offline-window consistency & general rewriting & windowed consistency proofs & no \\
G3 directory privacy & enumeration / recon & VRF-derived labels & no \\
G4 fork detection & A4 & cross-vantage auditor & yes ($\geq 1$ honest) \\
\bottomrule
\end{tabularx}
\end{table}

% ============================================================================
\FloatBarrier
\section{Out of Scope}\label{sec:threat-scope}

\begin{itemize}
  \item \textbf{Prevention.} A determined operator or network can always refuse
  service. The contribution is to ensure interference cannot masquerade as successful
  verification. A refusal or drop is observable as an unavailable or incomplete check,
  though its cause is not attributed.
  \item \textbf{Endpoint security.} A compromised device (X1) is outside the
  model. No directory mechanism can help once the adversary holds the victim's
  keys.
  \item \textbf{Query privacy from the operator.} The server learns which
  identities it is asked to resolve. DUET does not provide query privacy from the operator, a
  separate property from the directory privacy of Goal~G3 (\cref{sec:threat-goals}). A
  distributed-OPRF extension is noted as future work (\cref{ch:discussion}).
  \item \textbf{Group messaging.} The design is pairwise. Extending warning
  propagation to group settings such as Sender Keys or MLS is left to future work.
  \item \textbf{Multi-device keychains.} Extending the warning to the multi-device
  keychains ELEKTRA targets (\cref{sec:bg-vkd}) is out of scope. The model assumes one
  device per identity.
\end{itemize}

Each boundary is revisited, with the corresponding future-work direction, in
\cref{ch:discussion}.
